\documentclass{beamer}
\usetheme{Montpellier}
\DeclareOptionBeamer{compress}{\beamer@compresstrue}
\ProcessOptionsBeamer

\mode<presentation>

\useoutertheme[footline=authortitle]{miniframes}

\usecolortheme{dolphin}
\usecolortheme{orchid}

\definecolor{beamer@blendedblue}{rgb}{0.137,0.466,0.741}

\setbeamercolor{structure}{fg=beamer@blendedblue}
\setbeamercolor{titlelike}{parent=structure}
\setbeamercolor{frametitle}{fg=black}
\setbeamercolor{title}{fg=black}
\setbeamercolor{item}{fg=black}

\mode
<all>
\usepackage{appendixnumberbeamer}
\usepackage{pifont}

\setlength\fboxsep{-0.5pt}

\beamertemplatenavigationsymbolsempty
\usepackage{amsmath, amsthm, mathtools, amssymb, enumerate, physics, bm, bbm}
\newcommand{\E}{\mathbb{E}}
\newcommand{\R}{\mathbb{R}}
\newcommand{\N}{\mathbb{N}}
\newcommand{\foram}{~\forall~}
\newcommand{\forat}{\forall~}
\newcommand{\exist}{\exists~}
\newcommand{\sts}{~|~}
\newcommand{\stt}{~\text{such that}~}
\newcommand{\proofline}[2]{{#1}~\big(\text{#2}\big)}
\newcommand{\undersett}[2]{\underset{#1}{\text{#2}}~}
\newcommand{\parti}[2]{\frac{\partial #1}{\partial #2}}
\newcommand{\partitwo}[2]{\frac{\partial^2 #1}{\partial #2^2}}
\newcommand{\darti}[2]{\frac{d #1}{d #2}}
\newcommand{\dartitwo}[2]{\frac{d^2 #1}{d #2^2}}
\newcommand{\casesline}[2]{#1~\text{if}~#2}
\DeclareMathOperator*{\argmax}{arg\,max}
\DeclareMathOperator*{\argmin}{arg\,min}
\newcommand{\maxx}[2]{\max_{#1}#2}
\newcommand{\argmaxx}[2]{\argmax_{#1}#2}
\newcommand{\pars}[1]{\left(#1\right)}
\newcommand{\brac}[1]{\left[#1\right]}
\newcommand{\set}[1]{\left\{#1\right\}}
\newcommand{\bigp}[1]{\big(#1\big)}
\newcommand{\bigb}[1]{\big[#1\big]}
\newcommand{\bigs}[1]{\big\{#1\big\}}
\newcommand{\Bigp}[1]{\Big(#1\Big)}
\newcommand{\Bigb}[1]{\Big[#1\Big]}
\newcommand{\Bigs}[1]{\Big\{#1\Big\}}
\newcommand{\setof}[1]{\{#1\}}
\newcommand{\sumone}[2]{\sum_{#1 = 1}^{#2}}
\newcommand{\indic}{\mathbbm{1}}
\frenchspacing
\usepackage{tikz}
\usetikzlibrary{arrows.meta}
\usetikzlibrary{positioning}
\usetikzlibrary{calc}
\tikzset{%
  >={Latex[width=2mm,length=2mm]},
  % Specifications for style of nodes:
            base/.style = {rectangle, rounded corners, draw=black,
                           minimum width=2cm, minimum height=0.5cm,
                           text centered, font=\sffamily},
  PreMarket/.style = {base, fill=green!30},
  Normal/.style = {base, fill=blue!30},
  StrategyMethod/.style = {base, fill=red!30},
}

\title{The Impact of Environmental Taxation on CO2 Emissions}
\author{Sonia Mendizábal and Wanchao Zhao}
\date{Statistical Methods for Causal Inference}

\begin{document}
%\setbeamercovered{transparent}
\expandafter\def\expandafter\insertshorttitle\expandafter{%
  \insertshorttitle\hfill%
  \insertframenumber\,/\,\inserttotalframenumber}
%\addtobeamertemplate{navigation symbols}{}{ \hspace{1em}    \usebeamerfont{footline}%
%    \insertframenumber / \inserttotalframenumber }

\begin{frame}[noframenumbering]
\maketitle
\vspace{-10pt}
\begin{figure}
    \centering
    \includegraphics[width=6cm]{figures/vu_logo.png}
\end{figure}
\end{frame}

\section{Introduction}

\begin{frame}{Objective}

\end{frame}

\begin{frame}{Data}

\end{frame}


\section{Design}

\begin{frame}{Difference in Differences}

\end{frame}


\section{Theory}

\begin{frame}{Methodologies}

\end{frame}

\section{Results}

\begin{frame}{Methodology 1}

\end{frame}


\begin{frame}{Comparison}

\end{frame}

\section{Conclusions}

\begin{frame}{Continuity Assumption: Math}

\end{frame}



\section{References}

\begin{frame}{Continuity Assumption: Math}

\end{frame}

\end{document}
